\chapter{主理想整环上的自由模}

\begin{theorem}{左遗传环上自由模的子模结构}{}
	设环$A$的每个左理想作为左$R$-模都投射,$M$是自由$A$-模,则对$M$的每个子模$N$都同构于$A$的一族左理想的直和.
\end{theorem}

\begin{corollary}{左遗传环上投射模的子模依然投射}{}
	设环$A$的每个左理想作为左$A$-模都投射,$M$是投射左$A$-模,则$M$的每个子模都是投射模.
\end{corollary}

\begin{corollary}{PID上投射模=自由模}{}
	设$A$是PID,则每个投射$A$-模都是自由$A$-模.
\end{corollary}

\begin{remark}{}{}
	由此可知,对每个集合$I$,在$A$中有一族理想$\left(\mathfrak{a}_{i}\right)_{i\in I}$使得$A^{\left(I\right)}$和这族理想的直和同构.
\end{remark}

\begin{lemma}{交换环上有限维有限生成自由模的每个子模都有限维且自由}{}
	设$A$是交换环,$A$-模$M$由$n$个元素生成,则$M$的每个自由子模的维数都$\leq n$.
\end{lemma}

\begin{proposition}{PID上有限生成模的自由子模维数有限且自由}{}
	设$A$是PID.如果$M$是$n$维自由$A$-模,则$M$的每个子模都是维数$\leq n$的自由$A$-模.
\end{proposition}

\begin{corollary}{PID上有限生成模的子模的生成元个数有限}{}
	设$A$是PID,$A$-模$M$由$n$个元素生成,则$M$的每个子模最多由$n$个元素生成.
\end{corollary}